% Generated draft: requires mathematical and source review.
\documentclass[11pt]{article}
\usepackage[margin=25mm]{geometry}
\usepackage{fontspec}
\setmainfont{Latin Modern Roman}
\usepackage{amsmath,amssymb,mathtools,unicode-math}
\setmathfont{Latin Modern Math}
\usepackage{enumitem,needspace,xurl,hyperref}
\setlength{\parindent}{0pt}
\setlength{\parskip}{5pt}
\setlength{\emergencystretch}{3em}
\newcommand{\N}{\mathbb{N}}
\newcommand{\Z}{\mathbb{Z}}
\newcommand{\Q}{\mathbb{Q}}
\newcommand{\R}{\mathbb{R}}
\newcommand{\C}{\mathbb{C}}
\newcommand{\F}{\mathbb{F}}
\newcommand{\A}{\mathbb{A}}
\newcommand{\PP}{\mathbb{P}}
\newcommand{\E}{\mathbb{E}}
\newcommand{\eps}{\varepsilon}
\newcommand{\dd}{\,\mathrm{d}}
\begin{document}
\section*{UnsolvedMath ID: 5300042}
\textbf{Source problem code:} catalogue entry
\subsection{Definitions and mathematical statement}
Let $f:\widehat\C\to\widehat\C$ be a rational map of degree at least $2$, the usual degree range for Julia-set dynamics. Its Julia set $J(f)$ is the complement of the open set on which the family of iterates is normal. Write
\[
 \omega_f(z)=\bigcap_{N\geq0}\overline{\{f^n(z):n\geq N\}}
\]
for the forward omega-limit set, with closure in the sphere. Let $m$ denote spherical area, normalized to total mass one; this gives the usual Lebesgue-null-set convention on the sphere.

The rational map is nonsingular for this measure class, but area preservation is not assumed. A standard conservativity condition is that every measurable $A$ with $m(A)>0$ has
\[
 m(A\cap f^{-n}A)>0\quad\text{for some integer }n\geq1.
\]
This is a nonsingular recurrence property, not an assertion that $m$ is an invariant probability.

\textbf{Questions (Mikhail Lyubich).} For every rational $f$ satisfying $J(f)=\widehat\C$, is it true that
\[
 m\{z:\omega_f(z)=\widehat\C\}=1,
\]
and is $f$ conservative with respect to Lebesgue/spherical area? The first question concerns almost-everywhere full forward accumulation; the second concerns recurrence of arbitrary positive-area measurable sets. Neither is restricted to an unstated special family of rational maps.

The supplied research summary marks both conclusions solved positively and attributes them to Lyubich's ergodic theory and Eremenko--Lyubich, calling them classical and saying they were verified through the literature. However, no article, theorem number, publication details, or exact hypotheses are provided. This draft records that claimed status without independently adopting the verification claim. The background separately says the source lists the item open and no comprehensive modern resolution check was available, so exact all-rational-map coverage requires source review.
\subsection{Short English statement}
If a rational map has Julia set equal to the whole sphere, must almost every point have the whole sphere as its omega-limit set, and must the map be conservative for spherical area? The registry claims positive solutions but supplies no precise theorem authenticating their full scope.
\subsection{Sources}
\begin{itemize}
\item Local UnsolvedMath registry, permanent ID 5300042, source code catalogue entry, ``Conservativity when the Julia set is the sphere''; proposer: Mikhail Lyubich, proposed year 1992. Supplied status solved, classification SOLVED-IN-LITERATURE, updated 2026-08-06.
\item Stony Brook, \emph{Open Problems in Dynamical Systems}, 1992, Lyubich P2, \url{https://www.math.stonybrook.edu/open-problems-dynamical-systems}. Reported source-TeX-and-PDF-text extraction and access: 2026-07-29. The record says source-side partial results were preserved but no comprehensive modern resolution check was available. Rights: NEEDS\_REVIEW; public preprint access/source provenance verified but redistribution permission not established.
\item The summary attributes positive conclusions to Lyubich and Eremenko--Lyubich but gives no specific paper or theorem. Its statement of verification is itself supplied evidence, not work performed for this draft.
\end{itemize}
\end{document}
